%% kapa: exam=1 round=30 year=2019 period=1 form=A subject=civil_law subject_ko=민법 seq=36 number=36
%% source: Q-Net 감정평가사 기출문제 (cst00309) artlSeq=5205205 file=extracted/1008/5205205/2019년 제30회 감정평가사 1차시험 문제지/2019년 제30회 감정평가사 1차시험 문제지/2019년도 제30회 감정평가사 1차 1교시 A형.hwp
%% conversion: kordoc 4.24.51 -> tools/gen1.py (auto); review=proofread-sonnet; flags=none
\begin{question}{36}
전세권에 관한 설명으로 옳지 않은 것은?
\begin{choices}
  \item 전세권은 저당권의 목적이 될 수 있다.
  \item 전세권자와 인지(隣地)소유자 사이에도 상린관계에 관한 규정이 준용된다.
  \item 전세권자는 필요비 및 유익비의 상환을 청구할 수 있다.
  \item 전세권의 존속기간은 10년을 넘지 못한다.
  \item 전세금의 지급이 전세권의 성립요소이기는 하지만, 기존의 채권으로 전세금의 지급에 갈음할 수도 있다.
\end{choices}
\end{question}
\answer{3}
